% !TEX root = ../main.tex

\chapter{논의}
\label{ch:discussion}

본 장은 \ref{ch:results}장 결과를 RQ1--RQ4 및 SRQ1에 대해 해석하고, DeFi 리스크 평가에 대한 이론적·실무적 함의를 논의하며, 타당도 위협 요인을 평가하고 주요 EndorseRank--AWP 비교의 한계를 명시한다.

\dissertationSubheading[sec:findings-rq]{연구 질문별 결과}

BigQuery Tier~1 결과를 주요 연구 질문(RQ1--RQ4)별로 요약한다:

\textbf{RQ1 (평판 구조):} EndorseRank와 AWP는 부분적으로 겹치면서 구성개념적으로 구별된 순서를 산출한다(inter-method Kendall $\tau = 0.316$, 약~중간 구간). Allowance는 승인 의도를, 송금은 경제적 흐름을 부호화한다; 더 희소한 $G_{\text{approve}}$ 그래프는 $G_{\text{transfer}}$와 다른 중심성을 산출한다.

\textbf{RQ2 (사회적 방법 정렬):} EndorseRank와 AWP 중 구성개념별 우세 방법이 나타난다: allowance에서 EndorseRank($\tau=0.355$ vs.\ $-0.030$); 송금($0.534$ vs.\ $0.333$) 및 대부분 사회-결과 계열에서 AWP(표~\ref{tab:alignment-three-methods}). 거래 성공 프록시에서 구체적으로 AWP가 EndorseRank를 초과($0.179$ vs.\ $0.096$; 표~\ref{tab:alignment-gmx}). Allowance 기반 전파는 본 코호트에서 수익성 있는 무기한 청산과 송금 기반 AWP보다 더 강하게 정렬되지 않는다.

\textbf{RQ3 (효율):} EndorseRank의 AWP 대비 PageRank 런타임 이점(표~\ref{tab:benchmark-runtime}; $\approx 2.1$\,s vs.\ $\approx 18.7$\,s, Tier~1에서 $\approx 8.9\times$)은 감쇠와 함께 전체 송금 이력을 집계하는 대신 최신 allowance 스냅샷을 사용하기에서 비롯된다. Tier~2 풀 스케일링(표~\ref{tab:benchmark-scaling})은 $n=10{,}000$($\approx 10.4\times$) 및 $N=31{,}612$($\approx 9.4\times$)에서 격차를 확인한다. Damping 및 상위 토큰 강건성 검사(표~\ref{tab:robustness-damping}--\ref{tab:robustness-tokens})는 allowance--송금 트레이드오프 패턴을 바꾸지 않는다.

\textbf{RQ4 (검증 프레임):} 송금 프록시는 두 사회적 방법 모두에서 거래 성공보다 더 높은 평균 $\tau$를 산출한다(표~\ref{tab:alignment-three-methods}). 결과 기반 프록시는 사회적 그래프 전반에서 전반적으로 약~중간 정렬을 보여, 본 표본에서 단일 프록시 계열이 ``평판''을 완전히 포착하지 못함을 시사한다.

\dissertationSubheading[sec:supplementary-finding-srq1]{보조 결과 (SRQ1)}

GF-PR은 거래 성공 평균 $\tau$($0.583$)와 역리스크 평균 $\tau$($0.376$; 표~\ref{tab:alignment-three-methods})에서 최고치를 달성한다. 이는 해당 프록시 정의에 사용된 동일 PnL 스트림에서 구축된 결과-네이티브 그래프에 대해 예상된다. 이 상한은 거래 결과 계열에서 높은 정렬이 우수한 사회적 그래프 의미론이 아닌 구성개념 중첩에서 비롯될 수 있음을 명확히 한다. EndorseRank와 AWP는 배포 가능한 사회적 평판 방법으로 남고, GF-PR은 진단 참조만(주장 SC1).

\dissertationSubheading[sec:theoretical-implications]{이론적 함의}

5-프록시 프레임워크는 Do et al.의 검증 논리를 확장하여 송금 중심성, allowance 위임, 레버리지 결과, Sybil 조정 활동을 분리한다. EndorseRank는 주요 allowance 구성개념($G_{\text{approve}}$)과 정렬; 송금 프록시는 2차 검증 축(중간 $\tau=0.333$, AWP $0.534$ 미만). AWP는 송금---주요 그래프 구성개념--- 및 동일 지갑 활동에서 파생된 여러 결과 프록시와 정렬. 보조 GF-PR 분석은 사회적 엣지가 아닌 결과 흐름에 직접 순위를 매길 때 달성 가능한 거래 결과 프록시 정렬 상한을 보여 준다(SRQ1).

\dissertationSubheading[sec:practical-implications]{실무적 함의}

송금 허브 탐지가 필요한 프로토콜은 비용이 높더라도 AWP를 선호할 수 있다. 명시적 지출 승인이나 저지연 갱신을 강조하는 프로토콜은 EndorseRank를 선호할 수 있으며, 더 낮은 송금 프록시 정렬을 수용한다. GF-PR은 평가 설계 중 결과 프록시 중첩에 대한 진단 상한으로만 적합하며, 배포 가능한 사회적 평판 점수가 아니다. 공개 파이프라인은 allowance 로그가 있는 모든 ERC-20 체인에서 교차 프로토콜 재현을 가능하게 한다.

\dissertationSubheading[sec:threats-validity]{타당도 위협 요인}

실증 주장 해석 시 다음 위협을 고려해야 한다:

\begin{itemize}
\item \textbf{Proxy not ground truth:} 무담보 대출 부도 레이블 없음; 무기한 거래 청산 및 PnL(GMX V2)은 도메인 직관 프록시이며 대출 상환이 아님.
\item \textbf{GF-PR construct overlap:} GF-PR은 거래 성공 및 역리스크 프록시와 입력 의미론을 공유; 해당 계열의 높은 $\tau$는 우수한 사회적 평판으로 해석해서는 안 됨(SRQ1).
\item \textbf{Sybil and wash trading:} 적대적 allowance 또는 송금 패턴이 순위를 조작할 수 있음; Sybil 조정 프록시는 부분 완화만 제공.
\item \textbf{Single protocol and chain:} 다른 무기한 또는 대출 프로토콜로의 일반화는 교차 프로토콜 검증 필요.
\end{itemize}

\dissertationSubheading[sec:limitations]{한계}

결과는 Arbitrum One의 6개월 창을 반영한다. 토큰 소수점 정규화, USD 가격 오라클, permit 대 approve 커버리지가 엣지 가중에 영향을 줄 수 있다. 확장 baseline(LP-PR, CW-AWP, LF-PR, RiskProp, Aave W$\leftrightarrow$W 변형)은 향후 비교를 위해 보관(부록~\ref{ch:appendix-eval})하나 본문 서술에서는 제외.

\textbf{Scalability scope.} Tier~2 스케일링(표~\ref{tab:benchmark-scaling})은 $N=31{,}612$ 지갑 풀(GMX + 평판 서브그래프 parquet에서 파생한 보충 주소)을 사용한다. $n=10{,}000$ 및 전체 풀 크기 단계를 측정했으며, 제안 목표인 $50{,}000$/$100{,}000$ 지갑은 Tier~2 BigQuery 평판 재추출 완료를 기다린다(\texttt{extract\_reputation\_data.py --tier2}; dry-run $\approx 1.35$\,TiB 총량). Tier~1 in-cohort 스케일링($n \leq 5{,}521$)은 부록~\ref{ch:appendix-eval}에 유지.
